\documentclass{article}
\usepackage{amsmath,amssymb}
\usepackage[margin=1in]{geometry}

\newcommand{\KL}{\mathrm{KL}}
\newcommand{\Hred}{H_{\mathrm{red}}}
\newcommand{\Hfree}{H_{\mathrm{free}}}
\newcommand{\eps}{\epsilon}

\begin{document}

\title{Entropy of a full-support CCD}
\author{Yang, Bouckaert, Drummond}
\date{}
\maketitle

\noindent
We estimate the phylogenetic entropy of the full-support CCD regulariser of the
companion note. The standard CCD entropy is the wrong quantity here; we give two
estimators that agree with brute force where it is computable, apply them to the
RSV2 posterior (129 taxa), and use them to track how the escape parameter $\mu$
should change with sample size.

\section*{Why the standard CCD entropy fails}

For a plain CCD every tree probability is a product of conditional clade
probabilities (CCPs) $\theta$ over a finite set of observed clades, and the
entropy follows the Lewis recursion
$H=-\sum_C P(C)\sum_{(L,R)}\theta(L,R\mid C)\log\theta(L,R\mid C)$.
The full-support model breaks this assumption: at a reservable clade the
conditional distribution is a \emph{mixture} of an observed (red) split with
probability $1-\mu$ and an escape into a \emph{blue region} (novel clades) with
probability $\mu$, and the observed CCPs sum to $1-\mu$, not $1$. The Lewis sum
ignores the $(1-\mu)$ discount, the escape mass, and the entire novel-clade
support, so it returns the entropy of the renormalised observed backbone --- a
systematic underestimate. We therefore do not use the standard formula, and
estimate the entropy directly.

\section*{Two estimators}

We write $P(T)$ for the model probability of a tree: each red node contributes
$(1-\mu)\,\theta(S_v\mid C_v)$ and each maximal blue region with top clade $C$ and
$m-2$ novel clades contributes $\eps(C)^{m-2}$ (FLAT weighting; SHARED divides by
the boundary's pathcount). Here $\eps(C)$ solves
$\sum_{j\ge1}N_j(C)\,\eps^{j}=\mu$ over the computed reserve orders $j\le k$, with
$N_j$ the number of order-$j$ boundaries.

\paragraph{Sampling estimator.}
Entropy is the expected surprisal, so
\begin{equation}
  \hat H = -\frac{1}{S}\sum_{s=1}^{S}\log P(T_s),\qquad T_s\sim\text{model},
  \qquad \mathrm{SE}=\mathrm{sd}\!\big(\log P\big)\big/\sqrt{S}.
\end{equation}
The trees $T_s$ are drawn from the model itself and each carries its probability
$P(T_s)$, so the estimate is unbiased; its only cost is the $O(1/\sqrt{S})$
sampling noise.

\paragraph{Recursive estimator (generalised Lewis recursion).}
The model is generated by a recursive branching process over the observed-clade
DAG, so entropy decomposes by the chain rule. Let $\Hfree(C)$ be the entropy of
the subtree generated when $C$ may escape, and $\Hred(C)$ when $C$ is forced to
take an observed split (a blue region's boundary part). Then
\begin{align}
  \Hred(C) &= \sum_{(L,R)}\theta(L,R\mid C)\Big[-\log\theta(L,R\mid C)
              +\Hfree(L)+\Hfree(R)\Big],\\[2pt]
  \Hfree(C) &= \underbrace{(1-\mu)\big[-\log(1-\mu)+\Hred(C)\big]}_{\text{red branch}}
   \;\underbrace{-\,\log\eps(C)\sum_{m}(m-2)\,N_m\,\eps(C)^{m-2}}_{\text{blue self-information}} \notag\\
   &\qquad +\;\underbrace{\sum_{\text{regions}}\!P(\text{region})\!\sum_i \Hred(B_i)}_{\text{blue children}} ,
\end{align}
with $\Hfree(C)=\Hred(C)$ when $C$ reserves nothing. The red part and the blue
self-information are closed-form; the blue-children term is summed exactly by
enumerating the same canonical boundaries the reserve uses (FLAT weights each by
its pathcount). Computing $\Hfree,\Hred$ for clades in increasing size makes every
child and boundary part already available, so the recursion is deterministic and
\emph{exact} for the self-consistent distribution (escape mass $\mu$, regions up
to depth $k$).

\paragraph{Which distribution.}
The entropy reported is that of whichever distribution one chooses to draw from.
The \emph{self-consistent} distribution (escape mass exactly $\mu$, blue regions
only up to reserve depth $k$) is normalised and has the closed-form recursion
above. The \emph{full-support} distribution additionally carries the reserve tail
(regions deeper than $k$); this tail is the same intractable quantity that can
only be estimated by sampling, so the full-support entropy is obtained from the
sampling estimator rather than the recursion. The two agree only when the chosen
distribution and the tail correction are matched: no tail with the self-consistent
distribution, the sampled tail with the full-support distribution.

\section*{Validation against brute force}

On four taxa with no tail correction the model normalises exactly, so summing
$P(T)$ over all rooted topologies gives the exact entropy. The recursion matches
it to $10^{-9}$ and the sampling estimate to within its standard error, for both
weightings.

\begin{center}
\begin{tabular}{llrrr}
\hline
mode & $\mu$ & exact $H$ & recursive & sampling ($\pm$SE)\\
\hline
FLAT   & 0.05 & 1.033235 & 1.033235 & $1.0307\pm0.0017$\\
SHARED & 0.05 & 1.023676 & 1.023676 & $1.0269\pm0.0017$\\
FLAT   & 0.01 & 0.784942 & 0.784942 & $0.7835\pm0.0011$\\
\hline
\end{tabular}
\end{center}

\section*{RSV2 (129 taxa)}

The topology space is far too large to enumerate, so the two estimators
cross-validate each other. At $\mu=0.01$ (FLAT, $k=2$) they agree, and the
escape/novel-clade support adds about $11\%$ over the observed-only backbone that
the standard formula would report.

\begin{center}
\begin{tabular}{rrrrr}
\hline
trees & clades & backbone (obs.\ only) & recursive & sampling ($\pm$SE)\\
\hline
1{,}000  & 1143 & 51.23 & 59.19 & $59.02\pm0.11$\\
5{,}000  & 1630 & 52.79 & 59.31 & $59.17\pm0.08$\\
18{,}001 & 2009 & 53.31 & 59.10 & $59.08\pm0.08$\\
\hline
\end{tabular}
\end{center}

\noindent
On the full posterior the estimators agree to $0.24$ standard errors. (Entropies
in nats.)

\section*{Selecting $\mu$: dependence on sample size}

The entropy grows with $\mu$, so a fixed default is arbitrary: $\mu$ should be the
value that maximises cross-validated held-out tree log-probability. On RSV2 this
optimum is essentially independent of the backbone smoothing $\alpha$ (the same
grid point wins for every $\alpha$), so we fix $\alpha$ and grid $\mu$. Taking
disjoint, consecutive post-burn-in blocks of growing size, $\mu^{*}$ falls
steadily: a richer backbone observes more clades, held-out trees contain fewer
genuinely novel clades, and less escape mass is justified.

\begin{center}
\begin{tabular}{rrrrrr}
\hline
block & clades & $\mu^{*}$ & held-out (nat/tree) & $P(\text{novel})$ & $H(\mu^{*})$\\
\hline
128   & 730  & 0.02020 & $-62.42$ & 0.80  & 60.78\\
512   & 926  & 0.00553 & $-56.40$ & 0.34  & 56.04\\
2048  & 1278 & 0.00196 & $-55.37$ & 0.13  & 55.57\\
8192  & 1743 & 0.00061 & $-54.40$ & 0.041 & 54.78\\
\hline
\end{tabular}
\end{center}

\noindent
Over the $64\times$ span $\mu^{*}$ drops by $\sim33\times$, i.e.\
$\mu^{*}\propto N^{-0.8}$, a little slower than $1/N$; $P(\text{novel})$ collapses
in step. All four optima are interior. At these small $\mu^{*}$ the reserve tail
($\sim\mu^{3}$) is negligible, so the self-consistent recursion and the
full-support entropy coincide.

\section*{Held-out probability is a third entropy estimate}

The held-out average log-probability is, by construction, the negative
\emph{cross-entropy} between the true posterior $Q$ and the fitted model $P$:
\begin{equation}
  -\,\overline{\log P}_{\text{held-out}} = H(Q,P) = H(Q) + \KL(Q\Vert P).
\end{equation}
The reported $H(\mu^{*})$ is the model's \emph{self-entropy} $H(P)$; the two
differ only in whether surprisal is averaged under $Q$ or under $P$, so they
coincide when $P\approx Q$ --- which is exactly the regime $\mu^{*}$ targets, since
maximising held-out log-probability minimises $\KL(Q\Vert P)$ ($H(Q)$ being
independent of $\mu$). This is why the last two columns above nearly match, and it
gives three estimates of the same posterior entropy ($\approx 55$ nats):
the model self-entropy, the held-out cross-entropy (an upper bound), and brute
force on small trees.

The residual gaps $+1.64,+0.36,-0.20,-0.38$ are interpretable. At small $N$ the
model fitted on a training subset carries a real $\KL>0$, lifting the
cross-entropy above $H(P)$. As $N$ grows $\KL\to0$ and a small \emph{negative} gap
remains: the regularised model is genuinely flatter (higher entropy) than the
posterior. This is not a comparison artifact --- measuring the held-out
cross-entropy against the \emph{same} cross-validation models' self-entropy
(rather than the whole-block model's) gives essentially the same gap,
$\approx-0.44$ nats at $N=8192$. Two effects make the model flatter, in very
unequal measure: (i)~it over-reserves novel-clade mass --- model
$P(\text{novel})=0.041$ against a held-out rate of $0.034$, a $20\%$ excess --- but
the binary entropy of that mismatch is only $\approx0.02$ nats; (ii)~conditional on
a novel clade, the model spreads its mass nearly uniformly over an astronomically
large set of resolutions (high conditional entropy), whereas the posterior
concentrates on the few novel clades actually visited, which accounts for the
remaining $\approx0.43$ nats. The flatness is therefore mostly honest agnosticism
about \emph{which} unseen tree, not miscalibration of the \emph{total} novel mass
--- the intended behaviour of a full-support prior.

\paragraph{Weighting and sample size.}
The figures above use the FLAT (clustering-neutral) weighting. The SHARED weighting
shares a boundary's escape mass over its pathcount resolutions, penalising
\emph{clustered} novel clades. On the 8192-tree block SHARED is indeed the better
fit, but only marginally: held-out log-probability higher by $0.0033$ nat/tree, and
entropy at $\mu^{*}$ essentially unchanged ($54.783$ vs FLAT's $54.785$). The reason
is that at large $N$ the held-out-optimal $\mu^{*}\approx6\times10^{-4}$ is small
enough that escape is almost entirely \emph{single} novel clades --- the order at
which FLAT and SHARED coincide --- so the two weightings, which differ only on
clustered novelty (mass $\sim\eps^{2}$), are near-identical here. The distinction
bites at smaller $N$ (larger $\mu^{*}$); at the large-$N$ operating points the
entropy figures reported here are weighting-independent to three decimals.

\end{document}
